\documentclass[11pt,a4paper]{article}
\usepackage[utf8]{inputenc}
\usepackage[T1]{fontenc}
\usepackage[portuguese]{babel}
\usepackage{xcolor}
\usepackage{hyperref}
\usepackage{geometry}
\usepackage{tcolorbox}

\geometry{margin=2.5cm}

\definecolor{primary}{RGB}{39, 174, 96}
\definecolor{secondary}{RGB}{44, 62, 80}
\definecolor{codebg}{RGB}{240, 240, 240}

\hypersetup{colorlinks=true, linkcolor=primary, urlcolor=primary}

\title{\color{secondary}\Huge\textbf{Solução: Exercício 5.2 - Primeiros passos com a malha de serviço Istio}}
\author{\color{primary}DevOps com Kubernetes -- 2026}
\date{}

\begin{document}
\maketitle

\section*{\color{primary}Guia de Solução}
\begin{tcolorbox}[colback=green!5!white,colframe=green!60!black,title=Resolução Passo a Passo]
### Passo a Passo

1. \textbf{Instalação do CLI}: Siga o guia de documentação e baixe os binários do Istio. Ex: \texttt{curl -L https://istio.io/downloadIstio | sh -} e inclua o executável no seu PATH.
2. \textbf{Configurar k3d/Kubernetes}: Execute \texttt{istioctl install --set profile=demo -y} no cluster para instalar a estrutura do \textit{Control Plane} do Istio.
3. \textbf{Injeção Automática}: Rotule o \textit{namespace} padrão para usar o proxy automático do Istio: \texttt{kubectl label namespace default istio-injection=enabled} (ou as configurações específicas para a versão Ambient, sem sidecars).
4. \textbf{Deploy do Sample App}: Faça o apply dos manifestos fornecidos pelo repositório do Istio (como o \textit{Bookinfo application}).
5. \textbf{Monitoramento e Limpeza}: Exponha o \textit{Kiali} via \texttt{istioctl dashboard kiali}. Visualize os nós se comunicando. Em seguida, limpe as instalações conforme instruído pela documentação.
\end{tcolorbox}

\end{document}
